\scanpage{9}
\mspage{6}

ressemble d'assez près à $\mathcal X_2$ (les surfaces sous-jacentes sont homéom., \unsure{tout} comme les surfaces sous-jacentes $\mathcal X_0,\mathcal X_1,\mathcal X'_1$ le sont entre elles). On l'obtient en \deleted{\unsure{ajoutant}}\inserted{admettant} les positions relatives super\deleted{singulières}\inserted{critiques} parmi les faces admises, mais en \deleted{prenant \illegible{1}}\inserted{considérant}, pour une $D_i$ fixée ayant $n_i$ sommets, $2n_i$ positions relatives supercritiques \inserted{$D'_i$} correspondantes, suivant la position du pt d'intersection virtuel $D_i\cap D'_i$, qui doit être sur une des $2n_i$ facettes sur $\unsure{\widetilde D_i|\widehat D_i}$.

On trouve encore que tous les sommets sont d'ordre $4$. La \deleted{\unsure{ps.dr.}} ps. droite $D_i\subset\mathcal X_2$ est ici remplacée par une \deleted{\unsure{collier}\ \illegible{1}}\inserted{bande (de Möbius)} de $D_i$. Cela revient, à partir de $\mathcal X'_1$, \deleted{au lieu} \deleted{alors}\inserted{après} avoir découpé les faces supercritiques, au lieu d'identifier directement des pts antipodiques de chaque comp. conn. du bord, \deleted{de}\inserted{à} \deleted{prendre}\inserted{\unsure{fixer}} d'abord \deleted{\unsure{une}} une ``virole'' à ce bord \unsure{par un} bord, qui délimite cette bande de Möbius.

Cette fois on aura une application $\mathcal X_3\mathbin{-}X$, telle que l'image inverse de $X_1$ soit ce qu'il faut, \deleted{Mais} y compris les $D_i$ de \unsure{tout à l'h.} Mais on notera que sur les faces \inserted{supercritiques} de $\mathcal X_3$, il n'y a pas cette fois-ci de représentation affinisée de $\underline\Sigma=(X,\Sigma)$ comme dans le cas de $\mathcal X'_1$, \deleted{mais} plutôt une représentation d'une partie pertinente d'un voisinage tubulaire conv. de $D_i$ dans $X$.

D'autre part, les \deleted{\unsure{7}} faces supersingulières ne sont pas de $=$ \deleted{ordre que} les autres, de sorte que $\mathcal X_3$ ne se prête pas à une interprétation, en termes de la combinatoire interne, du syst. local des $\operatorname{Pol}(D)$.

\margin{\uline{TSVP}}
\margin{8}
